\documentclass[10pt,a4paper]{amsart}
\usepackage[margin=19mm]{geometry}
\usepackage{amsmath,amssymb,mathtools}
\usepackage[scr=rsfs,cal=euler]{mathalfa}
\usepackage{xfrac}
\usepackage[unicode,hidelinks,bookmarks=false]{hyperref}
\newcommand{\ie}{i.e.\ }
\newcommand{\cf}{cf.\ }
\newcommand{\resp}{respectively\ }
\newcommand{\Subsection}{\subsection}
\providecommand{\nobreakdash}{\nobreak\hbox{-}}
\setlength{\emergencystretch}{2em}
\multlinegap0pt
\usepackage{amssymb}
\usepackage{dsfont}
\usepackage{xcolor}
\newtheorem{proposition}{Proposition}
\newtheorem{theorem}{Theorem}
\title{Reflected Brownian Motion in the half plane with mixed boundary reflexions}
\author{Jules Flin}
\date{Source: arXiv:2604.27951v2 (CC BY 4.0)}
\begin{document}
\maketitle

\noindent\emph{Source and licence.} Reflected Brownian Motion in the half plane with mixed boundary reflexions. Version arXiv:2604.27951v2 (CC BY 4.0), distributed by arXiv under the Creative Commons Attribution 4.0 International licence, http://creativecommons.org/licenses/by/4.0/, as recorded in the arXiv licence metadata for this exact version and shown on the abstract page https://arxiv.org/abs/2604.27951v2. Full licence notice: \texttt{licenses/arxiv-2604.27951-CC-BY-4.0.txt}.\par\smallskip

\noindent\textit{Editorial scope.} Counted unit: the article's proposition prop:ex (source0.tex lines 454--458). The article's complete original proof of it is delivered with it (source0.tex lines 460--487, one \texttt{proof} environment of 28 lines). No mathematics is written, repaired or re-derived: the statement and the proof are the article's own lines, in the article's own notation, and no retained line is substituted or re-flowed.  Retained uncounted as the source-local context the proof needs: lines 150--152; lines 366--373; lines 417--434.  Nothing was omitted from any retained span, so the package carries no omission record.  No retained line contains a citation, so the wrapper carries no bibliography and no bibliography entry.  No retained line contains a figure environment, an image-inclusion command or drawing code, so no image is delivered and the package declares no asset dependency.  Editorial formatting only, in addition: an AMS \texttt{amsart} preamble that restates the article's own declarations and macros used by the retained lines, namely the environments \texttt{proposition}, \texttt{theorem} on separate counters, together with the standard packages those lines need.  The retained environments print in this document as Proposition 1, Theorem 1 in order of appearance, whereas the article prints the same items with the numbering of its own full document.  The complete native source of the article as distributed in the arXiv e-print for this exact version is bundled unchanged as \texttt{sources/199512/source0.tex}.
\begin{equation}\label{eq:FE}
    K(x,y)\phi(x,y)+\sum_\pm k_\pm(x,y)\phi_\pm(x)=0,
\end{equation}

\begin{proposition}[Sokhotski-Plemelj formula]\label{prop:SP} Let $f\in \breve{H}^\mu$ for some $\mu\in (0,1]$. The function
$$F(z):=\frac{1}{2i\pi}\int_{\mathbb R}\frac{ f(\tau)}{\tau-z}\mathrm{d}\tau,\quad z\in\mathbb C\setminus \mathbb R,$$
is sectionally analytic, and admits, for all $t\in\mathbb R$, a limit from below $F_-(t)$ and a limit from above $F_+(t)$, that satisfy
\begin{equation}\label{eq:jump_cond}
    F_+(t)-F_-(t)=f(t),\quad F_+(t)+F_-(t)=\frac{1}{i\pi }\;\mathrm{p.v.}\int_{\mathbb R}\frac{f(\tau)}{\tau-t}\mathrm{d}\tau.
\end{equation}
Moreover, $F$ is bounded at infinity.
\end{proposition}

\begin{theorem}[Explicit Laplace transforms]\label{thm:laplace} 
The Laplace transforms $\phi_\pm(x)$ are given by the following explicit formulas:
\begin{enumerate}
    \item[(a)] For $\pm \mathfrak{Re}(x) < 0$,
    \begin{equation*}
        \phi_\pm(x) = \frac{\Lambda}{(x \mp 1)^{1-\alpha}} \exp\left(\frac{1}{2i\pi} \int_{\mathbb{R}} \frac{\log\widetilde{G}(\tau)}{\tau+ix} \mathrm{d}\tau \right).
    \end{equation*}
    \item[(b)] On the imaginary axis, for $x=it$ with $t \in \mathbb{R}$,
    \begin{equation*}
        \phi_\pm(it) = \frac{\Lambda}{(it \mp 1)^{1-\alpha}} \widetilde{G}(t)^{\pm 1/2} \exp\left(\frac{1}{2i\pi} \,\mathrm{p.v.} \int_{\mathbb{R}} \frac{\log\widetilde{G}(\tau)}{\tau-t} \mathrm{d}\tau \right),
    \end{equation*}
    where the principal value integral is defined as in Proposition~\ref{prop:SP}.
\end{enumerate}
Here, the common normalization constant $\Lambda$ is uniquely determined by the relation
\begin{equation}\label{eq:normalization}
    \phi_\pm(0) = \pm \frac{\mu_1 - \mu_2 r_\mp}{r_- - r_+}.
\end{equation}
\end{theorem}

\begin{proposition}[A symmetric case]\label{prop:ex} If $(\mu_1,\mu_2,r_-,r_+)=(0,-1,1,-1)$, then the density $\pi$ of the invariant measure is given by
$$\pi(u,v)=\frac{1}{\sqrt{\pi}}\frac{\sqrt{\sqrt{u^2+v^2}+v}}{\sqrt{u^2+v^2}}\exp\left(-\sqrt{u^2+v^2}-v\right),$$
or equivalently, in polar coordinates
$$\pi(r\cos(\theta),r\sin(\theta))=\sqrt{\frac{2}{\pi r}}\cos\left(\frac{\theta}{2}-\frac{\pi}{4}\right)\exp\left(-2r\cos^2\left(\frac{\theta}{2}-\frac{\pi}{4}\right)\right).$$
\end{proposition}

\begin{proof}
    With this choice of parameters, the function $\widetilde{G}(t)$ evaluates to a constant (in fact, it is the only choice that leads to such simplification). Hence, Theorem~\ref{thm:laplace} yields 
    \begin{equation}\label{eq:lateral_laplace}
        \phi_\pm(x)=\frac{1}{2\sqrt{1\mp x}}.
    \end{equation}
    We recognize here the Laplace transforms of a Gamma distribution, and the corresponding lateral densities are given explicitly by
    $$\pi_\pm(u)=\frac{1}{2\sqrt{\pi}}\frac{\exp(-|u|)}{\sqrt{|u|}}.$$
    Injecting~\eqref{eq:lateral_laplace} into the functional equation~\eqref{eq:FE} leads to the following expression for the bivariate Laplace transform:
    \begin{equation}\label{eq:laplace_symmetric}
    \phi(x,y)=\left(\frac{1}{\sqrt{1-x}}+\frac{1}{\sqrt{1+x}}\right)\frac{1}{1+\sqrt{1-x^2}-y}.
    \end{equation}
    First, we invert $\phi(x,y)$ with respect to the second variable, by noticing that, as a function of $y$, it is of the form $c/(y_0-y)$. Here $y_0= y_0(x)=1+\sqrt{1-x^2}$. Hence, we obtain the following \textit{partial inverse} identity
    \begin{equation*}
        \int_{-\infty}^{+\infty}\pi(u,v)e^{xu}\mathrm{d}u=\left(\frac{1}{\sqrt{1-x}}+\frac{1}{\sqrt{1+x}}\right)\exp\left(-v\left(1+\sqrt{1-x^2}\right)\right)
    \end{equation*}
    For $x=it$, this reduces our problem to the Fourier inversion of this quantity with respect to $t$:
    $$\pi(u,v)=\frac{e^{-v}}{2\pi}\int_{-\infty}^{+\infty}\left(\frac{1}{\sqrt{1-it}}+\frac{1}{\sqrt{1+it}}\right)\exp\left(-v\sqrt{1+t^2}-itu\right)\mathrm{d}t.$$
    Successively applying the substitutions $t=\sinh \xi$ and $\omega=\exp(\xi/2)$, the integral in the above expression reduces down to
    \begin{align*}
        &\quad\!\quad\int_{-\infty}^{+\infty}\left(\frac{1}{\sqrt{1-it}}+\frac{1}{\sqrt{1+it}}\right)\exp\left(-v\sqrt{1+t^2}-itu\right)\mathrm{d}t\\
        &=2\int_{-\infty}^{+\infty}\cosh\left(\frac{\xi}{2}\right)\exp\left(-v\cosh \xi -iu\sinh \xi\right)\mathrm{d}\xi\\
        &=2\int_0^{+\infty}\exp\left(-a\omega^2-\frac{\overline{a}}{\omega^2}\right)\mathrm{d}\omega+2\int_0^{+\infty}\frac{1}{\omega^2}\exp\left(-a\omega^2-\frac{\overline{a}}{\omega^2}\right)\mathrm{d}\omega\\
        &=2\int_0^{+\infty}\exp\left(-a\omega^2-\frac{\overline{a}}{\omega^2}\right)\mathrm{d}\omega+2\int_0^{+\infty}\exp\left(-\overline{a}\omega^2-\frac{a}{\omega^2}\right)\mathrm{d}\omega,
    \end{align*}
    where $a=\frac12(v+iu)$. Some standard hyperbolic trigonometric and algebraic simplifications are omitted. Here, we recognize Gaussian-type integrals, which yields
    $$\pi(u,v)=\frac{\exp(-v-2|a|)}{2\sqrt{\pi}}\left(\frac{1}{\sqrt{a}}+\frac{1}{\sqrt{\overline{a}}}\right).$$
    Minor algebraic cleanups (\textit{e.g.}, $2|a|=\sqrt{u^2+v^2}$) conclude the proof.
\end{proof}

\end{document}
